\scanpage{10}
\hfill\handwritten{24.4,67}\\

Dear Murre,\\

I am sending you enc\unsure{d|h|l}osed the sketch which I promised on ramificatio\unsure{n}
% The last cell of ramification is cut by the physical right edge: its left stem and top serif survive, but its right-hand part does not.
data, as well as your manuscript together with my annotations. As the paragraphs 1 to 5 will have to be changed entirely, there is maybe not much point for you to look up to carefully my annotations on these. \deleted{As} I would only suggest that you put more fo\unsure{rm|rmm}ulas, arrows etc on \deleted{sep\unsure{ara}te} separate lines, for easier reading; that you normalize the sta\unsure{\deleted{dn}nd|\deleted{nd}dn}ard references to SGA, EGA (for instance, no need to indicate SGA when reference is being made to the seminar 1960-61 in which your talk is going to be included itself, reference as (XI 6.1) is then sufficient); that you be cautious when using a locution like ``if primitive $n$.th roots exist over $S$'', which does not mean much by itself, unless you make clear what you really mean. I would propose the following tentative plan for your talk:

\begin{enumerate}[label={}]
\item[1.] Kummer coverings with respect to a given set of sections.
\item[2.] \reversemarginpar
\margin{\unsure{\handwritten{\figure{\begin{tikzpicture}
\draw[gray!45,line width=.25pt] (0.30,0.85) .. controls (0.37,0.59) and (0.37,0.27) .. (0.25,0.16);
\draw[gray!55,line width=.3pt,dash pattern=on 1pt off 1.2pt] (0.25,0.16) .. controls (0.18,0.07) and (0.06,0.02) .. (0,0);
\end{tikzpicture}}}}}
\normalmarginpar Ramification data.
% Faint grey curved mark in the left margin beside item 2, extending towards item 3. Its status as an intentional hand stroke rather than a paper feature is uncertain.
\item[3.] Properties of the catégory of coverings of type $\underline{R}$. The group ${}_{1}^{\underline{R}}(S,\ )$
% Underlined R is stacked directly over 1. No base glyph is visible; the second argument is blank. The closing parenthesis reaches the physical edge; no following point is visible.
\item[4.] Functorial behaviours.
\item[5.] Galois coverings of type $\underline{R}$.
\item[6.] Abhyankars theorem (for divisors with normal crossings).
\end{enumerate}
etc

The sections 6 to 10 essentially the same as in your notes, just taking into account my observations, and changing inner references according to the changes made in 1. to 5. A paragraph 11 could be an appendix on the tame variant of Lefschetz's theorem. - Par.1 is essentially your previous par.1 ; you should check exactly what is needed in it for \unsure{\deleted{f}w|\deleted{w}f}hat
% No terminal punctuation is visible on this page.
% archive sheet number: 9
